

%\section{You \emph{can} have an appendix here.}


\section{Training with MMCR for Generative Models} \label{t_mmcr}

\begin{figure} [t]
	\centering
	%\vspace{-2mm}
	\hspace{-2mm}
	\includegraphics[width=0.52\linewidth]{Contents/figures/Representation_with_MMCRs}
	\hspace{-4mm}
	\includegraphics[width=0.52\linewidth]{Contents/figures/Representation_w_o_MMCRs}
	\vspace{-8mm}
	\caption{An example of models learned representation with or without MMCRs on MNIST. MMCRs support learning a separate representation, inspiring us to unlearn samples focusing only on the corresponding manifold representation. %\vspace{-2mm} 
	} \vspace{-2mm}
	\label{representations_w_or_w_o_mmcrs}
\end{figure}

 
 \Cref{representations_w_or_w_o_mmcrs} demonstrates the comparison of training a self-supervise model with or without MMCRs \citep{yerxa2023learning} on MNIST. With MMCRs during model training, the representation will make every class's points stay close together (small intra-class distance) while different classes lie far apart (large inter-class distance). The representation with these properties inspires us to investigate an efficient approximate unlearning method within only the manifold representation. 
 
 % Requires: \usepackage{booktabs}
 
 \begin{table}[h]
 	\centering
 	\caption{Frequently used notations in ManiF-SMC. \vspace{-2mm}}
 	\label{tab:notation_manif_smc}
 	\footnotesize
 	\setlength{\tabcolsep}{6pt}
 	\renewcommand{\arraystretch}{1.0}
 	\begin{tabular}{@{}l p{0.385\textwidth}@{}}
 		\toprule
 		Symbol & Meaning \\
 		\midrule
 
 		$\theta_o$ & Original model parameters (trained on $S$).\\
 		$\theta_u$ & Unlearned model parameters after forgetting $S_u$.\\
 		$\tilde{\theta}$ & Surrogate model sampled on the SMC path, $\tilde{\theta}=\phi_w(t^\star)$.\\
 
 		$S$ & Original training set.\\
 		$S_u$ & Unlearning (erased) set.\\
 		$S_r$ & Remaining (retained) set, $S_r = S\setminus S_u$.\\
 		$x_i$ & An erased target sample, $x_i\in S_u$.\\
 		$S_k^i$ & Top-$k$ most similar retained samples for $x_i$, selected in the original representation space; $S_k^i\subseteq S_r$.\\
 		$S_k$ & Union of retained neighborhoods, $S_k=\cup_{x_i\in S_u} S_k^i$.\\
 		$k$ & Number of retained neighbors used for each erased sample.\\
 		$\texttt{z}_{i,o}$ & Original representation of $x_i$: $\texttt{z}_{i,o}=f_{\theta_o}(x_i)$.\\
 		$\texttt{z}_{i,u}$ & Unlearned representation of $x_i$: $\texttt{z}_{i,u}=f_{\theta_u}(x_i)$.\\
 		$\texttt{c}_{i,o}$ & Neighbor centroid under the original model (computed on $S_k^i$ using $f_{\theta_o}$).\\
 		$\texttt{c}_{i,u}$ & Neighbor centroid under the unlearned model: $\texttt{c}_{i,u}=\frac{1}{|S_k^i|}\sum_{x_j\in S_k^i} f_{\theta_u}(x_j)$.\\
 		$\texttt{c}_{i,\tilde{\theta}}$ & Neighbor centroid estimated by the surrogate model: $\texttt{c}_{i,\tilde{\theta}}=\frac{1}{|S_k^i|}\sum_{x_j\in S_k^i} f_{\tilde{\theta}}(x_j)$.\\
 		
 
 		$w$ & Learnable control point of the quadratic B\'ezier path.\\
 		$t$ & Path parameter, sampled from $\mathcal{U}[0,1]$ when optimizing $w$.\\
 		$t^\star$ & Fixed sampling position on the path (e.g., $t^\star=0.5$) used to form $\tilde{\theta}$.\\
 
 
 		\bottomrule
 	\end{tabular}
 	\vspace{-2mm}
 \end{table}
 
 \section{Notations in ManiF-SMC} \label{notations_in_ManiF_SMC}
 
 We summarize the notations in \Cref{tab:notation_manif_smc}.
 
 
 \begin{algorithm}[h]
 	\caption{ManiF-SMC} \label{algManiF_SMC}
 	\begin{normalsize}
 		\BlankLine
 		\KwIn{Trained model $\theta_o$; unlearning set $S_u$; retained neighbor sets $\{S_k^i\}_{x_i\in S_u}$ and $S_k=\cup_{x_i\in S_u} S_k^i$; distance $\mathrm{dist}(\cdot,\cdot)$; control point $w$; learning rate $\eta$; epochs $E$; fixed $t^\star$ (e.g., $0.5$).}
 		\KwOut{Updated model parameters $\theta_u$.}
 		
 		\BlankLine
 		\tcp{Initialization}
 		$\theta_u \gets \theta_o$ \tcp*{start from original model}
 		$w \gets \theta$ \tcp*{initialize Bézier control point}
 		\ForEach{$x_i \in S_u$}{
 			$\texttt{z}_{i,o} \gets f_{\theta_o}(x_i)$ \tcp*{cache original representation}
 		}
 		
 		\BlankLine
 		\For{$e=1$ \KwTo $E$ epochs}{
 			\tcp{(A) Learn SMC control point $w$ on retained neighborhood data}
 			\ForEach{mini-batch $B_k \subset S_k$}{
 				sample $t \sim \mathcal{U}[0,1]$\;
 				$\theta_t \gets \phi_w(t) = (1-t)^2\theta_u + 2t(1-t)w + t^2\theta_o$\;
 				$\mathcal{L}_{\mathrm{path}} \gets \mathcal{L}_{B_k}(\theta_t)$ \tcp*{retained loss (same as training loss)}
 				$w \gets w - \eta \nabla_w \mathcal{L}_{\mathrm{path}}$\;
 			}
 			
 			\BlankLine
 			\tcp{(B) Pick a surrogate model on the low-loss path}
 			$\tilde{\theta} \gets \phi_w(t^\star)$ \tcp*{we pick $t^\star=0.5$ in experiments} %\;
 			
 			\BlankLine
 			\tcp{(C) ManiF update: compute centroid + adaptive margin under $\tilde{\theta}$, update $\theta_u$}
 			\ForEach{mini-batch $B_u \subset S_u$}{
 				\ForEach{$x_i \in B_u$}{
 					$\texttt{c}_{i,\tilde{\theta}} \gets \frac{1}{|S_k^i|}\sum_{x_j \in S_k^i} f_{\tilde{\theta}}(x_j)$\;
 					$\alpha_i \gets \Big[\mathrm{dist}\big(f_{\tilde{\theta}}(x_i),\texttt{z}_{i,o}\big) - \mathrm{dist}\big(f_{\tilde{\theta}}(x_i),\texttt{c}_{i,\tilde{\theta}}\big)\Big]_+$\;
 					$\ell_i \gets \Big[\mathrm{dist}\big(f_{\theta_u}(x_i),\texttt{c}_{i,\tilde{\theta}}\big) - \mathrm{dist}\big(f_{\theta_u}(x_i),\texttt{z}_{i,o}\big) + \alpha_i\Big]_+$\;
 				}
 				$\mathcal{L}_{\mathrm{triplet}} \gets \sum_{x_i \in B_u} \ell_i$\;
 				$\theta_u \gets \theta_u - \eta \nabla_{\theta_u}\mathcal{L}_{\mathrm{triplet}}$\;
 			}
 		}
 		\BlankLine
 		\Return $\theta_u$\;
 	\end{normalsize}
 \end{algorithm}
 

 
 
\section{Implementation Details of ManiF-SMC and Discussion} \label{algorithm_appendix}
 

 
\textbf{Overview.}
Algorithm~\ref{algManiF_SMC} summarizes ManiF-SMC, which alternates between (i) constructing a retained-geometry surrogate via self mode connectivity (SMC) and (ii) updating the unlearning model parameters via the manifold contrastive forgetting (ManiF) objective. The key design is that both the retained-neighbor centroid and the triplet margin are computed under a surrogate model sampled on a low-loss path trained only on retained neighborhood data, avoiding the need for retraining-from-scratch on $S_r$.

\paragraph{Inputs and preprocessing.}
The algorithm takes the trained model $\theta_o$, the erased (unlearning) set $S_u$, and for each $x_i\in S_u$ a retained neighborhood $S_k^i$ (top-$k$ most similar samples from the remaining data), with $S_k=\cup_{x_i\in S_u} S_k^i$. In practice, $S_k^i$ can be built once using the original representation space (e.g., nearest neighbors of $\texttt{z}_{i,o}=f_{\theta_o}(x_i)$) and then kept fixed during unlearning to reduce overhead.
We also cache $\texttt{z}_{i,o}$ for all erased samples (Lines 3--4), since it is used as the ``negative'' reference in the triplet constraint and does not change throughout the process.

We should notice that this algorithm does not rely on the label information. Hence, our inputs do not include the $y_i$.

\paragraph{Step (A): learning the SMC path control point $w$.}
Given endpoints $\theta_u$ (current unlearning model) and $\theta_o$ (original model), we parameterize a quadratic B\'ezier curve
$\phi_w(t) = (1-t)^2\theta_u + 2t(1-t)w + t^2\theta_o$.
In Step (A), we update the control point $w$ by minimizing the retained loss along the path (Eq.~\eqref{eq:w_opt}) using only mini-batches from $S_k$ (Lines 5--10). This yields a low-loss connection in parameter space that approximates the retained-data geometry without requiring full retraining on $S_r$. Importantly, the gradient is taken \emph{w.r.t. $w$ only}; $\theta_u$ is held fixed during this inner update so that the path training remains lightweight.

\paragraph{Step (B): sampling a surrogate model $\tilde{\theta}$.}
After optimizing $w$ on the retained neighborhood, we select a surrogate model $\tilde{\theta}=\phi_w(t^\star)$ (Line 11). We use a fixed $t^\star$ (e.g., $0.5$) for simplicity and stability. Empirically, mid-path sampling often provides a good trade-off between staying close to the retained geometry and not overfitting to either endpoint, while preserving the ``low-loss'' property by construction.

\paragraph{Step (C): ManiF update with centroid and adaptive margin.}
For each erased sample $x_i$, we compute the retained-neighbor centroid under the surrogate model,
$\texttt{c}_{i,\tilde{\theta}}=\frac{1}{|S_k^i|}\sum_{x_j\in S_k^i} f_{\tilde{\theta}}(x_j)$,
and an adaptive margin
$\alpha_i=\big[\mathrm{dist}(f_{\tilde{\theta}}(x_i),\texttt{z}_{i,o})-\mathrm{dist}(f_{\tilde{\theta}}(x_i),\texttt{c}_{i,\tilde{\theta}})\big]_+$
(Lines 14--15), which corresponds to Eq.~\eqref{adaptive_margin}. We then update $\theta_u$ by minimizing the triplet hinge objective (Eq.~\eqref{eq_triplet_loss_with_function}) on mini-batches from $S_u$ (Lines 17--18). Conceptually, this enforces that the updated representation $f_{\theta_u}(x_i)$ moves at least $\alpha_i$ closer to the retained-neighbor centroid than to its original representation $\texttt{z}_{i,o}$, implementing the push--pull unlearning behavior in a single margin-ranking constraint.

\paragraph{Why the surrogate is necessary.}
A core difficulty in representation-based unlearning is that the desired centroid $\texttt{c}_{i,u}$ and a suitable margin $\alpha$ are naturally defined under a retrained model on $S_r$, which is unavailable in efficient approximate unlearning. ManiF-SMC addresses this by using SMC to construct $\tilde{\theta}$ that is trained only on retained neighborhood data. Since $\tilde{\theta}$ is sampled from a low-loss connecting path, it provides a stable estimate of retained-neighbor representations (supported by Proposition~\ref{prop:smc_logit2}) and yields a data-dependent margin $\alpha_i$ that adapts to local geometry rather than using a global constant.

\paragraph{Distance function choices.}
We use Euclidean distance by default for $\mathrm{dist}(\cdot,\cdot)$ in Algorithm~\ref{algManiF_SMC}, which is standard in metric learning. Alternative metrics (e.g., cosine distance and $L_2$ norm) can be plugged in without changing the algorithm structure. We provide some evaluations about distance metrics in \Cref{distance_metrics_eval}.

%In practice, Euclidean distance on $\ell_2$-normalized representations often improves scale stability, and cosine distance is a drop-in alternative when representation norms vary significantly.

\paragraph{Other hyperparameters.}
Neighborhood size $k$ and sampling $t$ of SMC will also influence the performance of \Cref{algManiF_SMC}. We also provide the additional evaluations in \Cref{add_top_k_ablation,t_eval}.


\paragraph{Complexity discussion.}
Compared with retraining on $S_r$, ManiF-SMC is efficient because it (i) computes centroids only over the retained neighborhoods $S_k^i$ rather than the full retained set, and (ii) optimizes the SMC path using only $S_k$. The dominant additional cost beyond the ManiF update is the forward/backward pass for path training of $w$ on $S_k$ and the centroid computation for each $x_i$ within a mini-batch. Both are controllable via $k$ and the number of SMC optimization steps.

 


 \begin{table}[h]
	\scriptsize
	\caption{Dataset statistics.} 
%	\vspace{-4mm}
	\label{dataset_table}
	\resizebox{\linewidth}{!}{
		\setlength\tabcolsep{7.5pt}
		\begin{tabular}{cccc}
			\toprule[0.8pt]
			Dataset & Feature Dimension  & \#. Classes & \#. Samples \\
			\midrule
			MNIST \citep{deng2012mnist} & 28×28×1 & 10 & 70,000  \\  
			\rowcolor{verylightgray}
			CIFAR10 \citep{krizhevsky2009learning} & 32×32×3 & 10 & 60,000  \\  
			CelebA~\citep{liu2018large} & 178×218×3 & 2 (Gender) & 202,599 \\
			\rowcolor{verylightgray}
			Tiny-ImageNet \citep{le2015tiny} & 64×64×3 & 200 & 110,000 \\
			\bottomrule[0.8pt]
	\end{tabular}}
\vspace{-2mm}
\end{table}

\section{Datasets} \label{datasets_appendix}


%\noindent
%\textbf{Datasets.}
The statistics of all datasets used in our experiments are listed in \Cref{dataset_table}. Both MNIST and CIFAR10 are used to train 10-class classification models. The experiment on CelebA is to identify the gender attributes of the face images. The task is a binary classification problem, different from the ones on MNIST and CIFAR10. The task of Tiny-ImageNet is a 200-class classification. These datasets offer a range of objective categories with varying levels of learning complexity. We also introduce them as below.

%MNIST and CIFAR10 are for 10-class classification tasks, and Tiny-ImageNet is for 200-class classification tasks. The experiment on CelebA is to identify the gender attributes of the face images, which is a binary classification task.


\begin{itemize}[itemsep=0pt, parsep=0pt, leftmargin=*]
	\item \textbf{MNIST~\citep{deng2012mnist}.} MNIST contains 60,000 handwritten digit images for the training and 10,000 handwritten digit images for the testing. All these black and white digits are size normalized, and centered in a fixed-size image with 28 × 28 pixels.
	\item \textbf{CIFAR10~\citep{krizhevsky2009learning}.} CIFAR10 dataset consists of 60,000 32x32 colour images in 10 classes, with 6,000 images per class. There are 50,000 training images and 10,000 test images.
	\item  \textbf{CelebA~\citep{liu2018large}.} CelebA is a large-scale face attributes dataset with more than 200,000 celebrity images, each with 40 attribute annotations.
	\item \textbf{Tiny-ImageNet~\citep{le2015tiny} .} The Tiny-ImageNet image size is 64x64 pixels and the dataset sizes are 100,000 training images across 200 classes; 10,000 test images.
\end{itemize}

 \begin{figure*}[t]
	\centering
	\hspace{-2mm}
	\subfloat{  \label{fig:celebamodelmiauss}  \rotatebox{90}{ \hspace{12mm}	\small{On CelebA} }
		\includegraphics[scale=0.3]{../../../../PycharmProjects/Manifold_unl/Exp_results/On_CelebA/MBU_uss/CelebA_model_MIA_uss}
	}	
	\hspace{-2mm}
	\subfloat{ 	 	\label{fig:celebamodelrauss}
		\includegraphics[scale=0.3]{../../../../PycharmProjects/Manifold_unl/Exp_results/On_CelebA/MBU_uss/CelebA_model_RA_uss}
	}
	\hspace{-2mm}
	\subfloat{  		\label{fig:celebamodeltauss}
		\includegraphics[scale=0.3]{../../../../PycharmProjects/Manifold_unl/Exp_results/On_CelebA/MBU_uss/CelebA_model_TA_uss}
	}
%	\vspace{-2mm}
	\caption{Additional Evaluations of impact about different $\it USS$ on various unlearning methods on CelebA.   } %\vspace{-2mm}
	\label{evaluation_of_uss_cifar10_celeba} 
	%	\vspace{-6mm}
\end{figure*}


\section{Evaluation Metrics} \label{metric_detail}
We provide the details of metrics as follows.

\begin{itemize}[itemsep=0pt, parsep=0pt, leftmargin=*] % \vspace{-4mm}
	\item \textbf{Membership inference attack (MIA).} We employ the membership inference attack (MIA) \citep{song2019privacy,yeom2018privacy} to gauge the effectiveness of unlearning. Concretely, we apply the MIA predictor to the unlearned model $f_{\theta_u}$ on the dataset $S_u$. The resulting success rate of MIA is calculated by how many instances in $S_u$ are correctly identified as non-training samples for $f_{\theta_u}$. A higher MIA success rate suggests that $f_{\theta_u}$ retains less information about $S_u$.
	\item  \textbf{Remaining accuracy (RA).} This refers to the accuracy of $f_{\theta_u}$ on the remaining dataset $S_r$, which reflects the fidelity of machine unlearning. The training data information should be preserved from original model to the unlearned model.
	\item \textbf{Testing accuracy (TA).} We report TA as an indicator of how well the unlearned model $f_{\theta_u}$ generalizes when evaluated on the test dataset.
		\item  \textbf{Remaining mean-squared error (R-MSE).} This refers to the reconstruction mean-squared error of generative models on the remaining dataset $S_r$, which reflects the fidelity of machine unlearning for generative models.
	\item \textbf{Testing mean-squared error (T-MSE).} We report T-MSE as an indicator of how well the unlearned generative model performs when evaluated on the test dataset.
	\item  \textbf{Running Time (RT).} This metric reflects the computational efficiency of a machine unlearning method. We obtain RT by tracking the per-batch training duration and multiplying by the total number of training epochs.
\end{itemize}







\section{Machine Unlearning Benchmarks} \label{revisiting_mu}

Machine unlearning seeks to remove or negate the influence of a subset of training data $\{x_i\}_{i \in S_u}$ on a trained model $f_{\theta}$ \citep{bourtoule2021machine,warnecke2024machine}. Although exact unlearning can be achieved through retraining a model using the remaining dataset, the associated computational costs have driven the more efficient solutions, the approximate unlearning methods \citep{shen2024labelagnostic,izzo2021approximate}. Classic approximate unlearning approaches include:

\begin{itemize}[itemsep=0pt, parsep=0pt, leftmargin=*]
	\item \textit{Gradient Ascent (\textbf{GA})} \citep{graves2021amnesiac,thudi2022unrolling}: GA reverses the model training on the erased samples $S_u$ by adding the corresponding gradients back to $\theta_o$, i.e., moving $\theta_o$ in the direction of increasing loss for data to be unlearned, where $\theta_o$ is the original trained model parameters.
	%\item \textit{Influence Unlearning (\textbf{IU})} \citep{izzo2021approximate,sekhari2021remember}: IU utilize the influence funtion method \citep{koh2017understanding} to estimate how each sample $x_i$ in $S_u$ influenced $\theta_o$ and then ``subtract'' that influence from the model parameters.
	\item \textit{Variational Bayesian Unlearning (\textbf{VBU})} \citep{nguyen2020variational,nguyen2022markov}: VBU is an approximate unlearning method based on variational Bayesian inference. In practice, a middle layer of original neural networks is used as the Bayesian layer to calculate the VBU loss according to \citep{nguyen2020variational} to achieve unlearning.
	\item  \textit{Representation Forgetting Unlearning \textbf{(RFU)}} \citep{wang2023machine}: RFU tries to unlearn a bottleneck representation by minimizing the mutual information between the representation and the erased samples, which provides an unlearning solution from an information theory perspective. %We set a middle layer of the original model as the representation layer and implement unlearning according to \citep{wang2023machine}.
	\item  \textit{Saliency Unlearning \textbf{(SalUn)}} \citep{fan2024salun}: SalUn introduces ``weight saliency'' to remove the influence of samples and classes for unlearning a model, improving effectiveness and efficiency.
\end{itemize}

From the manifold representation viewpoint, unlearning a sample $(x_i,y_i) \in S_u$ can be achieved by pushing its representation $\texttt{z}_i$ away from the center of its learned local manifold representation, for example, the class-manifold centroid $\texttt{c}_{y_i}$. One simple operation proxy is to maximize the distance $\| \texttt{z}_i - \texttt{c}_{y_i} \|$, ensuring that the unlearning sample no longer ``lives'' near its old class manifold centroid. %The intuition raises the following approximate unlearning reformulation. 

% in \Cref{problem_df}.


%It is important to note that more favorable metrics for an approximate unlearning method should \textbf{reduce its performance disparity with the gold-standard retrained model (Retrain); except RT.} A higher value is not necessarily better. 

 



\section{Additional Evaluations}


\subsection{Additional Efficiency Evaluations about \textit{USS}} \label{add_exp_eff_uss}



We put the additional unlearning effectiveness evaluation on CelebA in \Cref{evaluation_of_uss_cifar10_celeba}, which shows the data removal effect and model utility preservation of ManiF-SMC.


We put the additional efficiency evaluation results on CelebA, in \Cref{additional_efficiency_of_uss}. The running-time curves show a clear efficiency hierarchy: the exact Retrain baseline is always the slowest—often by two to three orders of magnitude—while the approximate unlearning methods (VBU, ManiF, and SalUn) finish far sooner. Among the approximations, VBU is consistently the fastest, SalUn incurs the highest overhead, and ManiF-SMC sits in between, mirroring the computational complexity built into their respective update rules. 

 \begin{figure}[t]
	\centering
	\subfloat[\small On CelebA]{ 		\label{fig_celebarunningtimecsasbar}
		\includegraphics[scale=0.32]{../../../../PycharmProjects/Manifold_unl/Exp_results/On_CelebA/MBU_uss/CelebA_running_time_cs_as_bar}
	}
	\vspace{-2mm}
	\caption{Efficiency of Running time. Although the RT of approximate unlearning methods (VBU, ManiF-SMC, and SalUn) increases as the \textit{USS} increases, it is still much more efficient than retraining. \vspace{-0mm} } %\vspace{-2mm}
	\label{additional_efficiency_of_uss} 
	\vspace{-4mm}
\end{figure}


% As the number of points to be scrubbed (USS) grows from 200 to 1200, each approximate method’s cost rises monotonically but moderately (roughly linear in the log–log plots), whereas Retrain’s time remains essentially flat and dwarfs the others.  



%Dataset scale amplifies absolute timings—moving from MNIST to CIFAR-10 adds about one order of magnitude, and CelebA adds another—yet the relative savings over Retrain stay substantial: even on CelebA with 1 200 removals, the slowest approximation (ManiF) is still two orders of magnitude quicker than full retraining.  Overall, the figure highlights that approximate unlearning scales much more gracefully with deletion set size and dataset complexity, validating its practical advantage when strict exactness is unnecessary.
 
 
%\Cref{additional_efficiency_of_uss} compares the running time of approximate unlearning methods (VBU and ManiF-SMC) with retraining, for randomly selected unlearning sets ranging from 100 to 200 samples on MNIST (\Cref{fig_mnistrunningtimecsasbar}) and CelebA (\Cref{fig_celebarunningtimecsasbar}). As the size of the unlearning set increases, two approximate unlearning methods require more computation; however, VBU and ManiF-SMC maintain a consistently lower running time than retraining.




\begin{figure*}[t]
	\centering
%	\vspace{-2mm}
	\subfloat{		\label{fig:mnistmodelmiaussfinetuned} \rotatebox{90}{ \hspace{14mm}	\scriptsize{On MNIST} }
		\includegraphics[scale=0.3]{../../../../PycharmProjects/Manifold_unl/Exp_results/On_MNIST/MBU_uss_finetuned/MNIST_model_MIA_uss_finetuned}
	}			
	\subfloat{ 		\label{fig:mnistmodelraussfinetuned}
		\includegraphics[scale=0.3]{../../../../PycharmProjects/Manifold_unl/Exp_results/On_MNIST/MBU_uss_finetuned/MNIST_model_RA_uss_finetuned}
	}
	\subfloat{  		\label{fig:mnistmodeltaussfinetuned}
		\includegraphics[scale=0.3]{../../../../PycharmProjects/Manifold_unl/Exp_results/On_MNIST/MBU_uss_finetuned/MNIST_model_TA_uss_finetuned}
	}
	\vspace{-2mm} \\
	\subfloat{	\label{fig:cifar10modelmiaussfinetuned} \rotatebox{90}{ \hspace{14mm}	\scriptsize{On CIFAR10} }
		\includegraphics[scale=0.3]{../../../../PycharmProjects/Manifold_unl/Exp_results/On_CIFAR10/MBU_uss_finetuned/CIFAR10_model_MIA_uss_finetuned}
	}			
	\subfloat{ 		\label{fig:cifar10modelraussfinetuned}
		\includegraphics[scale=0.3]{../../../../PycharmProjects/Manifold_unl/Exp_results/On_CIFAR10/MBU_uss_finetuned/CIFAR10_model_RA_uss_finetuned}
	}
	\subfloat{  		\label{fig:cifar10modeltaussfinetuned}
		\includegraphics[scale=0.3]{../../../../PycharmProjects/Manifold_unl/Exp_results/On_CIFAR10/MBU_uss_finetuned/CIFAR10_model_TA_uss_finetuned}
	}
%	\vspace{-2mm}
	\caption{Additional evaluations for ManiF-SMC with fine-tuning utilizing the remaining dataset. ManiF-SMC with fine-tuning using the remaining dataset can effectively mitigate the model utility (RA and TA) degradation during unlearning, meanwhile the unlearning effectiveness (MIA) has slight drops compared with ManiF-SMC without fine-tuning. \vspace{-0mm} } %\vspace{-2mm}
	\label{evaluation_of_mcr_r} 
%	\vspace{-2mm}
\end{figure*}

 
 
 


\subsection{Limitations and Additional Evaluation with Fine Tuning} \label{with_finetuning}

\noindent
\textbf{Limitations of ManiF-SMC.} Since ManiF-SMC unlearns models only using the model manifold representation without the data task information, ManiF-SMC suffers greater utility loss compared with SalUn and other approximate unlearning methods that utilize the remaining data to fine-tune the unlearned model, especially when the unlearning sample size is large. Mitigating this degradation is therefore an important avenue for future research. In this work, we introduce a simple fine-tuning step on the remaining dataset and assess its effect on model performance.

\noindent
\textbf{Setup.} We additionally test the improvement of fine tuning for ManiF-SMC on MNIST and CIFAR10. The \textit{USS} is set from 200 to 1200, where 1200 is around $2\%$ training data on MNIST and CIFAR10. Our other settings are the same as the previous test for ManiF-SMC. The corresponding results are presented in \Cref{evaluation_of_mcr_r}. 
 
\noindent
\textbf{Evaluation Results.} In \Cref{evaluation_of_mcr_r},  we observed that ManiF-SMC with fine-tuning using the remaining dataset can effectively mitigate the model utility (RA and TA) degradation during unlearning, meanwhile the unlearning effectiveness (MIA) has slight drops compared with ManiF-SMC without fine-tuning. We should notice that the original ManiF-SMC unlearns only based on the model learned manifold representation, not needing the task information. Adding the fine-tuning on the remaining dataset will need label information to calculate the original learning loss. It improves the model utility preservation but makes the ManiF-SMC similar to the existing unlearning methods, having not cut the reliance on the task information.

 

\subsection{Ablation Study: How do Distance Metrics Influence the Performance of ManiF-SMC?} \label{distance_metrics_eval}


\noindent
\textbf{Setup.} 
As we mentioned above, different distance metrics may influence the unlearning methods. 
Therefore, we conducted the experiments using other metrics, the cosine similarity (the Normalized Temperature-Scaled Cross Entropy (NT-Xent) used in SimCLR \cite{chen2020simple}). For a better adaption of cosine similarity (NT-Xent), we also conducted Cosine Similarity (NT-Xent) with Multi-positive samples (add an additional one), because our method is for unlearning, the negative samples are limited.




\noindent
\textbf{Results.}
In \Cref{Different_Distance_Metrics}, the distance using cosine similarity achieves similar model utility preservation as the $L_2$-norm. However, the cosine similarity achieves a lower MIA value for unlearning effect than $L_2$-norm. We infer when using similarity as the triplet distance, the similarity has a higher preference to optimize the (anchor, positive) pair, making unlearning more difficult than the $L_2$-norm. For Cosine Similarity (NT-Xent) with Multi-positive samples, it does increase the unlearning model utility but further reduces the unlearning effect, the MIA. Moreover, we conducted experiments with the label information for fine tuning: $L_2$-norm (positive with label). It increases the unlearned model utility too and shows the potential of fine tuning for ManiF-SMC. However, the biggest advantage of our method is still not relying on the task and corresponding gradient ascent.


\begin{table}[t]
	\caption{Evaluation of Different Distance Metrics for ManiF-SMC on MNIST. 
	}
	\vspace{-2mm}
	\label{Different_Distance_Metrics}
	\resizebox{\linewidth}{!}{
		\setlength\tabcolsep{5.5pt}
		\begin{tabular}{ccccc}
			\toprule[0.8pt]
			\toprule[0.8pt]
			Distance Metric & MIA (\%)  & RA (\%) & TA (\%) & RT (second) \\
			\midrule
			$L_2$ norm &62.00 	& 99.59 	& 99.15 	&1.482 \\  
			\rowcolor{verylightgray}
			$L_2$ norm (with label) & 58.00  & 99.62 & 99.20 & 1.598  \\   %	53.99%	%	%
			Cosine Similarity (NT-Xent) &60.00 	& 99.38 &	99.04  & 	2.082  \\  
			\rowcolor{verylightgray}
			\makecell[c]{Cosine Similarity (NT-Xent)\\ with Multi-positive}   &54.00 & 99.55 & 99.15& 	2.703 \\ 
			\bottomrule[0.8pt]
	\end{tabular}}
 
\end{table} 


 \begin{table*}[t]
	\caption{Additional evaluations of the value of chosen $k$ remaining similar samples on MNIST, CIFAR10, and Tiny-ImageNet. \vspace{-4mm}
	}
	\small
	\label{ablation_study_k_remaining_samples2}
	\resizebox{0.98\linewidth}{!}{
		\setlength\tabcolsep{7.5pt}
		\begin{tabular}{ccccccccccccc}
			\toprule[0.8pt]
			\toprule[0.8pt]
			\multirow{2}{*} { \makecell[c]{\textbf{$k$} Value} } & \multicolumn{3}{c}{\textbf{On MNIST}} & \multicolumn{3}{c}{\textbf{On CIFAR10}} & \multicolumn{3}{c}{\textbf{On Tiny-ImageNet}}  &  \multicolumn{3}{c}{\textbf{On CelebA}} \\
			\cmidrule(lr){2-4}
			\cmidrule(lr){5-7}
			\cmidrule(lr){8-10}
			\cmidrule(lr){11-13}
			& MIA (\%)  & RA (\%) & TA (\%) & MIA (\%)  & RA (\%) & TA (\%) & MIA (\%)  & RA (\%) & TA (\%) & MIA (\%)  & RA (\%) & TA (\%)\\
			\midrule
			5 & 62.00 & 99.59 & 99.15    & 59.00 & 99.04 & 81.75   & 54.50 & 81.19 & 56.88  & 54.50 & 96.34 & 95.93\\  
			\rowcolor{verylightgray}
			6 &61.50 & 99.58 & 99.15      & 59.50 & 99.04 & 81.73   & 54.50 & 80.10 & 56.06   & 55.00 & 96.34 & 95.93  \\  
			7 & 61.00 & 99.56 & 98.84    & 59.50 & 99.04 & 81.73   & 54.50 & 80.10 & 56.07   & 55.00 & 96.35 & 95.94  \\  
			\rowcolor{verylightgray}
			8 &62.00 & 99.58 & 99.15    & 59.50 &99.03 & 81.75   & 54.50 &81.20 & 56.89   & 54.00 & 96.47 & 96.17 \\ 
			9 &61.50 & 99.58 & 99.15     & 59.50 & 99.04 & 81.75   & 54.50 & 81.20 & 56.91  & 54.00 & 96.47 & 96.17\\  
			\rowcolor{verylightgray}
			10 &61.50 & 99.58 & 99.15    & 58.50 & 99.05 & 81.75  & 54.50 & 81.19 & 56.90   & 53.00 & 96.48 & 96.19  \\  
			\bottomrule[0.8pt]
	\end{tabular}}
	%	\begin{tabbing}
		%		\hspace{2mm} --: the unlearning method is not applicable.
		%	\end{tabbing}
	% 	\vspace{-2mm}
\end{table*}






\subsection{Additional Study: Influence of $k$ Chosen Remaining Similar Positive Samples} \label{add_top_k_ablation}


%

\noindent
\textbf{Setup.} 
In ManiF-SMC, we choose the top-k remaining most similar samples as the positive set $S_k$, and the anchor is also calculated by the center of the top-k samples. Hence, we test how different top-k values influence our ManiF-SMC method. We randomly select 200 samples for unlearning, and the experimental setting is the same as above. We present the results on MNIST, CIFAR10, TinyImageNet, and CelebA in \Cref{ablation_study_k_remaining_samples2}. 


\noindent
\textbf{Results.}
Intuitively, more similar samples in the remaining dataset chosen will increase the model utility preservation during unlearning, because the changes of erasing one sample on the class-centroid can be expressed by $\|x_u - \texttt{c} \| / k$. The results on CelebA in \Cref{ablation_study_k_remaining_samples2} clearly confirm this, higher RA and TA when $k$ value increases. RA increases from 96.34\% ($k=5$) to 96.48\% ($k=10$), and TA climbs from 95.93\% to 96.19\%. However, the unlearning effectiveness drops at the same time when $k$ increases, from 54.50\% ($k=5$) to 53.00\% ($k=10$). Hence, we can claim that a larger $k$ assists to maintain the model utility but mitigates the unlearning effectiveness, which is also proved by other experiments on other datasets in \Cref{ablation_study_k_remaining_samples2}.

 

 

\iffalse
\begin{table}[t]
	%	\vspace{-4mm}
	\caption{Influence of the value of chosen $k$ remaining similar samples on CelebA.  \vspace{-4mm}
	}
	%	\small
	\label{ablation_study_k_remaining_samples}
	\resizebox{0.85\linewidth}{!}{
		\setlength\tabcolsep{8.5pt}
		\begin{tabular}{cccc}
			\toprule[0.8pt]
			\toprule[0.8pt]
			$k$ Value & MIA (\%)  & RA (\%) & TA (\%) \\
			\midrule
			5   & 54.50 & 96.34 & 95.93\\  
			\rowcolor{verylightgray}
			6 & 55.00 & 96.34 & 95.93  \\  
			7 & 55.00 & 96.35 & 95.94  \\  
			\rowcolor{verylightgray}
			8   & 54.00 & 96.47 & 96.17 \\ 
			9    & 54.00 & 96.47 & 96.17\\  
			\rowcolor{verylightgray}
			10  & 53.00 & 96.48 & 96.19  \\  
			\bottomrule[0.8pt]
	\end{tabular}}
	%	\begin{tabbing}
		%		\hspace{2mm} --: the unlearning method is not applicable.
		%	\end{tabbing}
	\vspace{-4mm}
\end{table}
\fi 


\begin{table}[t]
	% \vspace{-4mm}
	\caption{Evaluation of checkpoints sampled on the connectivity curve on MNIST. \vspace{-2mm}}
	\label{tab_mnist_connectivity_curve}
	\resizebox{0.95\linewidth}{!}{
		\setlength\tabcolsep{10.5pt}
		\begin{tabular}{cccc}
			\toprule[0.8pt]
			\toprule[0.8pt]
			$t$ & MIA (\%) & RA (\%) & TA (\%) \\
			\midrule
			0.1 & 58.50 & 99.62 & 99.22 \\
			\rowcolor{verylightgray}
			0.3 & 61.00 & 99.63 & 99.21 \\
			0.5 & 62.00 & 99.59 & 99.15 \\
			\rowcolor{verylightgray}
			0.7 & 58.50 & 99.59 & 99.20 \\
			0.9 & 58.00 & 99.62 & 99.21 \\
			\bottomrule[0.8pt]
		\end{tabular}
	}
	%	\vspace{-4mm}
\end{table}


\subsection{Ablation Study: How does $t$ of SMC Influence ManiF-SMC?} \label{t_eval}

 
	
In \Cref{SCM_section}, \Cref{mc_two,eq:w_opt} are applied to train the mode connectivity model to generate the local manifold representation for the adaptive margin calculation. We can train a connectivity curve, such as, $t \in [0,1]$ with step equal to 0.1, i.e., the curve with checkpoints in $\{0, 0.1, 0.2, 0.3, 0.4, 0.5, 0.6, 0.7, 0.8, 0.9, 1\}$. 
In our experiments, we use the checkpoint of $t = 0.5$, which is best for the unlearning effect. We also provide the experiments of the curve in the following \Cref{tab_mnist_connectivity_curve} to demonstrate the results.



\begin{table}[t]
	% \vspace{-4mm}
	\caption{Evaluation of different learning models on MNIST (before and after unlearning). \vspace{-4mm}}
	\label{mnist_other_learning_methods}
	\resizebox{0.95\linewidth}{!}{
		\setlength\tabcolsep{6.5pt}
		\begin{tabular}{c|cc|cc|cc}
			\toprule[0.8pt]
			\toprule[0.8pt]
			\multirow{2}{*}{Methods} 
			& \multicolumn{2}{c|}{MIA (\%)} 
			& \multicolumn{2}{c|}{RA (\%)} 
			& \multicolumn{2}{c}{TA (\%)} \\
			\cmidrule(lr){2-3}\cmidrule(lr){4-5}\cmidrule(lr){6-7}
			& Before unl. & After unl. & Before & After & Before & After \\
			\midrule
			Information Bottleneck (IB)      & 51.99 & 61.00 & 99.49 & 99.47 & 98.96 & 98.91 \\
			\rowcolor{verylightgray}
			InfoNCE & 50.50 & 53.50 & 99.39 & 99.40 & 98.91 & 98.83 \\
			\bottomrule[0.8pt]
		\end{tabular}
	}
	%	\vspace{-4mm}
\end{table}

\begin{table}[t] 
	% \vspace{-4mm}
	\caption{Evaluation of unlearning for ViT on MNIST (before and after unlearning). \vspace{-2mm}}
	\label{tab_vit_unlearning1}
	\resizebox{0.98\linewidth}{!}{
		\setlength\tabcolsep{6.0pt}
		\begin{tabular}{c|cc|cc|cc}
			\toprule[0.8pt]
			\toprule[0.8pt]
			\multirow{2}{*}{Model}
			& \multicolumn{2}{c|}{MIA (\%)}
			& \multicolumn{2}{c|}{MSE on remaining set}
			& \multicolumn{2}{c}{MSE on test set} \\
			\cmidrule(lr){2-3}\cmidrule(lr){4-5}\cmidrule(lr){6-7}
			& Before unl. & After unl. & Before & After & Before & After \\
			\midrule
			ViT & 62.00 & 67.00 & 0.0439 & 0.0443 & 0.0442 & 0.0445 \\
			\bottomrule[0.8pt]
		\end{tabular}
	}
	%	\vspace{-4mm}
\end{table}



\subsection{Application: Scalability to Unlearn Other Models} \label{Scalability_unlearn}


ManiF-SMC unlearning method is also transferable to other representation learning frameworks. We provide the experimental results for InforNCE (contrastive learning) and information bottleneck (IB) (representation learning) in \Cref{mnist_other_learning_methods}.




We also conducted new experiments for Vision Transformer (ViT) on MNIST for the image generative task. Results are provided in \Cref{tab_vit_unlearning1}, indicating the unlearning effectiveness of our method.

